problem:  
Let $\Sigma^n$ ($n\ge7$) be an exotic sphere that admits a Riemannian metric with positive sectional curvature. Consider the **moduli space**  
\[
\mathcal{P}(\Sigma^n) = \{\, g \text{ Riemannian on }\Sigma^n \mid \sec_g>0,\ \mathrm{Vol}(g)=1 \,\}/\mathrm{Diff}(\Sigma^n),
\]  
the space of volume‑normalized positively curved metrics modulo diffeomorphisms.  
For each $g\in\mathcal{P}(\Sigma^n)$, let $g_t$ denote the normalized Ricci flow starting at $g$ (assuming short‑time existence).  

**Question:**  
Does the normalized Ricci flow on $\Sigma^n$ necessarily escape every compact subset of $\mathcal{P}(\Sigma^n)$?  
That is, must every positively curved metric on an exotic sphere evolve—before converging—to a regime where curvature degenerates (e.g. sectional pinching collapses, injectivity radius tends to zero, or volume concentration occurs) rather than remaining within a compact region of the moduli space?  

Equivalently, could one have **Ricci flow stability** within $\mathcal{P}(\Sigma^n)$ comparable to that of the round sphere, or is such stability topologically obstructed on exotic spheres?  
Formally: if a sequence of normalized Ricci flows $g_t^{(k)}$ on $\Sigma^n$ converges to a smooth compact limit flow in $\mathcal{P}(\Sigma^n)$, must $\Sigma^n$ be diffeomorphic to the standard sphere?  

Why is it a "good" problem:  
This problem reframes the exotic sphere curvature‑existence issue into a **dynamical compactness question** in the moduli space of positively curved metrics under Ricci flow.  
It merges three modern geometric pillars:  

1. **Geometric analysis:** investigating normalized Ricci flow trajectories, compactness, and curvature blow‑up mechanisms.  
2. **Riemannian moduli theory:** the structure of $\mathcal{P}(M)$ under curvature and volume normalizations.  
3. **Differential topology:** the role of smooth structures in shaping asymptotic flow dynamics.  

Its statement is simple—asking whether exotic smooth structures obstruct geometric compactness under curvature evolution—but its consequences would be profound: a positive answer would give a new analytic characterization of the standard smooth structure via flow stability.  

The problem is mathematically well‑posed, avoids degeneracy or trivial scaling, and offers a bridge between Ricci flow dynamics and high‑dimensional topology, embodying the clarity, depth, and cross‑field integration expected of an outstanding research‑level problem.